\section{来源审计与问题边界：什么叫 Mechanistic Interpretability}

本讲追问一个比“模型是否答对”更困难的问题：给定一组已经训练好的参数，能否把模型内部计算反编译成人类可以检查的算法？Chris Olah 将这种工作称为 \term{mechanistic interpretability}（机制可解释性）。它并不满足于展示相关 activation 或生成自然语言理由，而是希望说明某组 weights、features 与 attention heads 如何组成一个可复现的计算过程。

本次重写以官方视频 \texttt{pC4zRb\_5noQ} 为影音基准。Stanford CS25 V1 课程表把讲座列在 2021-11-15；课堂中 Olah 多次称 induction-head 结果尚未发表。后来文章在 2022-03-08 公开，Stanford Online 到 2022-07-17 才上传录像。讲义因此将课堂中的 tentative hypothesis、后来文章的补充证据与本讲的工程解释明确分层，不能用后见之明把每个说法写成定论。

视频没有公开独立 slide PDF。我们从官方 1080p 录像中恢复 64 张教学页，保留公式推导的独立步骤与最终 progressive state，省略重复曲线、空白转场、片尾，以及 50--55 分钟低清 Lexoscope live demo 的中间界面；demo 中的 cross-lingual/soft-induction 口头解释仍完整进入正文。

\lecturefigure{slide-01-interpretability-definition.jpg}{Mechanistic interpretability：把 neural-network parameters 映射回人类可理解的 algorithms。}{Stanford Online 官方视频 00:00:30--00:03:27。}

图左的 vision circuit 给出最小例子：wheel feature 喜欢图像底部的轮子、不喜欢顶部的轮子；window feature 具有相反的空间偏好；它们共同支持 car feature。重点不是“发现一个会叫 car 的 neuron”，而是看到 features 之间有方向、符号与位置关系，可以被解释成一段组合算法。

\teachervoice{Olah 的类比是：weights 像 compiled computer program，features/neurons 像 variables 或 registers。机制解释的任务是 reverse engineering，而不是为输出写一段听起来合理的事后故事。}

\begin{knowledgebox}{Interpretability 不是一个单一问题}
\begin{tabularx}{\textwidth}{L{0.24\textwidth}X X}
\toprule
层次 & 典型问题 & 本讲是否重点处理 \\
\midrule
Behavioral evaluation & 模型在什么输入上成功或失败？ & 用来发现 ICL phase change，但不是终点。 \\
Feature attribution & 哪些 token/activation 与输出相关？ & 提供线索，单独不能证明 circuit。 \\
Mechanistic analysis & 参数怎样实现可执行算法？ & 本讲核心：QK/OV、path composition、induction。 \\
Causal validation & 移除组件后，目标行为会变化吗？ & small-model ablation 是关键证据。 \\
System/social explanation & 数据、部署、人与制度如何共同产生结果？ & 超出本讲电路范围，不能由 circuits 自动回答。 \\
\bottomrule
\end{tabularx}
\end{knowledgebox}

\subsection{为什么要做这种逆向工程}

Olah 给出两类动机。科学动机是：人类不会手写一个能完成 ImageNet classification 或 GPT-3 behavior 的明确程序，但 gradient descent 能在参数中找到某种解；理解这个解可能发现新算法。安全动机是：如果系统被用于高影响场景，我们希望在 failure mode 被实际触发之前，识别参数中潜伏的“unknown unknowns”。

\lecturefigure{slide-02-why-mechanistic-interpretability.jpg}{为什么研究机制解释：发现 gradient descent 写出的算法，并提前识别未知失败模式。}{Stanford Online 官方视频 00:03:28--00:04:23。}

\begin{warningbox}{能解释一个 circuit，不等于解释整个模型}
局部电路可以在特定模型、特定 prompt family 和特定 metric 上被验证；它不自动覆盖 MLP、训练数据、其他 heads、分布外行为或社会部署后果。Mechanistic interpretability 的价值来自可检验的局部进展，而不是提前宣称“模型已经透明”。
\end{warningbox}

\subsection{课堂当时的证据强度}

这堂课的特殊之处在于讲者主动展示研究的未完成状态。他此前主要研究 vision ConvNets，把同样方法迁移到 language models 仍属新工作；因此我们需要保留“这是 hypothesis”“large model 目前只有 correlation”“可能有其他 mechanism”这些限定语。

\lecturefigure{slide-03-language-models-new-territory.jpg}{讲者的版本声明：语言模型机制解释仍很新，离成熟发表与完整解释有距离。}{Stanford Online 官方视频 00:04:24--00:05:39。}

\teachervoice{Olah 欢迎听众提出最基础的问题，也欢迎指出不清楚和错误之处。这个课堂语气非常重要：后面的 induction-head story 是一条有多种证据支持的研究路线，不是对所有 Transformer in-context learning 的最终封闭解释。}

\subsection{本章小结}

Mechanistic interpretability 把模型视为被 gradient descent 编译的程序，并尝试恢复其中的算法。它既服务于科学发现，也可能帮助安全审计，但必须区分局部 circuit、整体模型与部署系统，并严格记录每条结论的证据强度。

